% cnbook/latex/base.tex
% bookdown 中文项目（XeLaTeX + ctexbook）公共件。
\usepackage{booktabs}
\usepackage{longtable}
% pandoc 为带列宽的表格生成 >{\raggedright\arraybackslash}p{…} 列，
% \arraybackslash 由 array 宏包提供；列宽里的算术表达式由 calc 宏包提供
\usepackage{array}
\usepackage{calc}

% xcolor 需带 usenames,dvipsnames：bookdown 模板把这两个选项传给 color
% （\PassOptionsToPackage{usenames,dvipsnames}{color}），而 hyperref 2022
% 不再加载 color，改由这里的 xcolor 提供 Maroon 等颜色名（目录链接色）
\usepackage{framed}
\usepackage[usenames,dvipsnames]{xcolor}
\definecolor{shadecolor}{RGB}{248,248,248}

\renewcommand{\textfraction}{0.05}
\renewcommand{\topfraction}{0.8}
\renewcommand{\bottomfraction}{0.8}
\renewcommand{\floatpagefraction}{0.75}

\makeatletter
\newenvironment{kframe}{%
\medskip{}
\setlength{\fboxsep}{.8em}
 \def\at@end@of@kframe{}%
 \ifinner\ifhmode%
  \def\at@end@of@kframe{\end{minipage}}%
  \begin{minipage}{\columnwidth}%
 \fi\fi%
 \def\FrameCommand##1{\hskip\@totalleftmargin \hskip-\fboxsep
 \colorbox{shadecolor}{##1}\hskip-\fboxsep
     % There is no \\@totalrightmargin, so:
     \hskip-\linewidth \hskip-\@totalleftmargin \hskip\columnwidth}%
 \MakeFramed {\advance\hsize-\width
   \@totalleftmargin\z@ \linewidth\hsize
   \@setminipage}}%
 {\par\unskip\endMakeFramed%
 \at@end@of@kframe}
\makeatother

\makeatletter
\@ifundefined{Shaded}{
}{\renewenvironment{Shaded}{\begin{kframe}}{\end{kframe}}}
\@ifpackageloaded{fancyvrb}{%
  % https://github.com/CTeX-org/ctex-kit/issues/331
  \RecustomVerbatimEnvironment{Highlighting}{Verbatim}{commandchars=\\\{\},formatcom=\xeCJKVerbAddon}%
}{}
\makeatother

\usepackage{makeidx}
\makeindex

\urlstyle{tt}

% PDF 书签带上章节编号：hyperref 默认把 \numberline（章节号）从书签里
% 去掉，书签只剩标题；bookmarksnumbered 让书签与目录一样显示完整编号
\hypersetup{bookmarksnumbered=true}

\usepackage{amsthm}
\makeatletter
\def\thm@space@setup{%
  \thm@preskip=8pt plus 2pt minus 4pt
  \thm@postskip=\thm@preskip
}
\makeatother

% 中文主字体与粗体：优先用思源宋体（Noto Serif CJK SC，SIL OFL 自由许可、
% 跨平台），BoldFont 配其真 Bold 字重——思源全部字重共享同一字符集
% （对 hong 全书用字实测：各字重均只缺 5 个扩展 B 区字，且都已在项目的
% \CJKfallbackchar 回退清单内），不存在「粗体比正文少字」的问题。
% Scale=1.02：思源宋体同字号下字面比宋体 SC 小、不放大则阅读显小显累，
% 默认放大 1.02（经 A/B/C/D/E 样张对比调试；想更大可用 1.06/1.10，
% 想笔画更饱满可换 Noto Serif CJK SC Medium）。
% 无思源的环境回退到 macOS 宋体 + AutoFakeBold=3 合成粗体：ctex 的 macOS
% 字体集默认 BoldFont=Songti SC Bold 与 Regular 字重差太小，__粗体__
% （\textbf）在 PDF 里几乎看不出；注意不能用 Songti SC Black——该字重
% 字符覆盖严重不全（hong 全书用字实测缺 136 字，如第 78 回「姽婳」，
% 粗体标题与目录章条目里直接丢字）。都没有时保持 ctex 默认，不干预
\IfFontExistsTF{Noto Serif CJK SC}{%
  \IfFontExistsTF{Kaiti SC}{%
    \setCJKmainfont[Scale=1.02, BoldFont=Noto Serif CJK SC Bold, ItalicFont=Kaiti SC]{Noto Serif CJK SC}%
  }{%
    \setCJKmainfont[Scale=1.02, BoldFont=Noto Serif CJK SC Bold]{Noto Serif CJK SC}%
  }%
}{%
  \IfFontExistsTF{Songti SC Regular}{%
    \setCJKmainfont[AutoFakeBold=3, ItalicFont=Kaiti SC]{Songti SC Regular}%
  }{}%
}
